\begin{center}
    \textbf{BAB I}\\
    \textbf{PENDAHULUAN}
    \addcontentsline{toc}{chapter}{BAB I. PENDAHULUAN}
\end{center}

\textbf{A. Latar Belakang}
\addcontentsline{toc}{section}{A. Latar Belakang}

\hspace*{1cm}Matematika merupakan salah satu mata pelajaran yang berperan penting dalam mengembangkan kemampuan berpikir logis, kritis, analitis, dan sistematis peserta didik. Namun, dalam praktik pembelajaran di sekolah, matematika masih sering dianggap sebagai mata pelajaran yang sulit, khususnya pada materi geometri seperti bangun ruang sisi datar. Kesulitan ini disebabkan oleh sifat materi yang abstrak serta keterbatasan media pembelajaran yang digunakan. Safira dan Salsabila (2021) menyatakan bahwa pembelajaran bangun ruang sisi datar memerlukan media visual yang mampu membantu siswa memahami objek secara konkret agar konsep yang dipelajari lebih mudah dipahami.

\hspace*{1cm}Kesulitan siswa dalam memahami materi bangun ruang sisi datar juga didukung oleh hasil penelitian empiris. Chintia, Amelia, dan Fitriani (2021) menemukan bahwa siswa mengalami kesulitan dalam memahami konsep dasar bangun ruang, menentukan rumus, serta menyelesaikan soal yang berkaitan dengan geometri ruang. Hal ini menunjukkan bahwa kemampuan visualisasi siswa terhadap objek tiga dimensi masih rendah, sehingga berdampak pada pemahaman konsep matematika secara keseluruhan.

\hspace*{1cm}Penelitian lain yang dilakukan oleh Putra, Amelia, dan Fitriani (2022) menunjukkan bahwa banyak siswa melakukan kesalahan dalam menentukan luas permukaan dan volume bangun ruang sisi datar. Kesalahan tersebut terjadi karena siswa kesulitan menghubungkan konsep geometri dengan representasi visual bangun ruang. Selain itu, Syafi’ah, Rusdiana, dan Ikmawati (2022) menyatakan bahwa siswa mengalami kesulitan dalam memahami unsur-unsur bangun ruang, jaring-jaring, serta penyelesaian soal kontekstual yang berkaitan dengan bangun ruang sisi datar.

\hspace*{1cm}Kesulitan tersebut juga berkaitan dengan rendahnya kemampuan visualisasi spasial siswa. Berdasarkan teori Van Hiele, kemampuan berpikir geometri siswa terdiri dari beberapa tahap, dan sebagian besar siswa masih berada pada tahap visualisasi. Sahara dan Nurfauziah (2021) mengungkapkan bahwa sebagian besar siswa belum mampu mencapai tahap analisis dalam memahami geometri ruang, sehingga mengalami kesulitan dalam memahami hubungan antar unsur bangun ruang sisi datar.

\hspace*{1cm}Seiring dengan perkembangan teknologi, inovasi dalam media pembelajaran menjadi salah satu solusi untuk mengatasi permasalahan tersebut. Media pembelajaran berbasis simulasi 3D memungkinkan siswa untuk melihat objek bangun ruang secara lebih nyata dan interaktif. Hidayah et al. (2024) menyatakan bahwa penggunaan teknologi berbasis Unity 3D dan Augmented Reality dapat meningkatkan hasil belajar siswa karena mampu menyajikan visualisasi objek secara konkret.

\hspace*{1cm}Selain itu, penggunaan media berbasis teknologi juga dapat meningkatkan motivasi belajar siswa. Setiawan et al. (2023) menyatakan bahwa media pembelajaran berbasis Augmented Reality mampu meningkatkan minat dan motivasi belajar siswa karena memberikan pengalaman belajar yang lebih menarik dan interaktif. Dengan adanya visualisasi objek tiga dimensi, siswa dapat lebih aktif dalam proses pembelajaran.

\hspace*{1cm}Berdasarkan uraian tersebut, dapat disimpulkan bahwa diperlukan pengembangan media pembelajaran yang mampu membantu siswa memahami materi bangun ruang sisi datar secara lebih konkret dan interaktif. Oleh karena itu, penelitian ini berfokus pada pengembangan media pembelajaran matematika berbasis simulasi 3D sebagai upaya untuk meningkatkan pemahaman konsep, kemampuan visualisasi spasial, dan motivasi belajar siswa.

\vspace{0.5cm}

\textbf{B. Identifikasi Masalah}
\addcontentsline{toc}{section}{B. Identifikasi Masalah}

\hspace*{1cm}Berdasarkan uraian pada latar belakang, dapat diidentifikasi beberapa permasalahan sebagai berikut:
\begin{enumerate}
    \item Siswa mengalami kesulitan dalam memahami konsep bangun ruang sisi datar karena sifat materi yang abstrak dan membutuhkan kemampuan visualisasi spasial yang tinggi.
    \item Masih banyak peserta didik kelas VIII SMP Negeri 3 Kalikajar yang memperoleh nilai ulangan harian pada materi bangun ruang sisi datar di bawah Kriteria Ketuntasan Minimal (KKM).
    \item Belum tersedianya media pembelajaran berbasis simulasi 3D yang dikembangkan secara khusus untuk membantu siswa memahami materi bangun ruang sisi datar.
\end{enumerate}

\textbf{C. Pembatasan Masalah}
\addcontentsline{toc}{section}{C. Pembatasan Masalah}

\hspace*{1cm}Berdasarkan identifikasi masalah, penelitian ini difokuskan pada pengembangan media pembelajaran berbasis Augmented Reality (AR) pada materi bangun ruang sisi datar guna mempermudah peserta didik dalam memvisualisasikan objek geometri tiga dimensi. Penelitian ini bertujuan untuk menguji tingkat kelayakan media yang dikembangkan, tanpa mengukur pengaruh media tersebut terhadap hasil belajar peserta didik.\\

\textbf{D. Rumusan Masalah}
\addcontentsline{toc}{section}{D. Rumusan Masalah}

\hspace*{1cm}Berdasarkan latar belakang dan identifikasi masalah, maka rumusan masalah dalam penelitian ini adalah sebagai berikut:
\begin{enumerate}
    \item Bagaimana prosedur pengembangan media pembelajaran berbasis Augmented Reality pada materi bangun ruang sisi datar?
    \item Bagaimana tingkat kelayakan media pembelajaran berbasis Augmented Reality pada materi bangun ruang sisi datar?
\end{enumerate}

\textbf{E. Tujuan Pengembangan}
\addcontentsline{toc}{section}{E. Tujuan Pengembangan}

\hspace*{1cm}Berdasarkan rumusan masalah di atas, maka tujuan dari penelitian ini adalah sebagai berikut:
\begin{enumerate}
    \item Mengetahui prosedur pengembangan media pembelajaran berbasis Augmented Reality pada materi bangun ruang sisi datar.
    \item Mengetahui tingkat kelayakan media pembelajaran berbasis Augmented Reality pada materi bangun ruang sisi datar.
\end{enumerate}

\textbf{F. Spesifikasi Produk yang Dikembangkan} \\
\addcontentsline{toc}{section}{F. Spesifikasi Produk yang Dikembangkan}
\hspace*{1cm}Produk yang dikembangkan oleh peneliti memiliki spesifikasi sebagai berikut:
\begin{enumerate}
    \item Produk berupa website yang dapat diakses melalui ponsel maupun laptop secara online.
    \item Produk mengintegrasikan materi bangun ruang sisi datar dengan teknologi Augmented Reality.
    \item Produk menampilkan objek AR melalui pemindaian kode QR (barcode).
\end{enumerate}

\textbf{G. Manfaat Pengembangan}
\addcontentsline{toc}{section}{G. Manfaat Pengembangan}

\hspace*{1cm}Hasil dari penelitian ini diharapkan dapat memberikan manfaat dalam dunia pendidikan, baik secara teoritis maupun praktis, sebagai berikut:

\begin{enumerate}
    \item \textbf{Manfaat Teoretis} \\
    Penelitian ini diharapkan memberikan kontribusi pengetahuan dalam bidang pendidikan matematika, khususnya terkait pengembangan media pembelajaran berbasis Augmented Reality.

    \item \textbf{Manfaat Praktis}
    \begin{enumerate}
        \item \textbf{Bagi Peserta Didik} \\
        Media ini diharapkan dapat membantu peserta didik dalam memvisualisasikan objek tiga dimensi, memahami materi bangun ruang sisi datar dengan lebih baik, serta meningkatkan minat belajar matematika.
        
        \item \textbf{Bagi Pendidik} \\
        Media ini dapat menjadi alternatif dalam menyampaikan materi bangun ruang sisi datar dan motivasi bagi pendidik untuk mengembangkan media pembelajaran berbasis Augmented Reality pada materi lainnya.
    \end{enumerate}
\end{enumerate}

\textbf{H. Asumsi dan Keterbatasan Pengembangan}
\addcontentsline{toc}{section}{H. Asumsi dan Keterbatasan Pengembangan}

\begin{enumerate}
    \item \textbf{Asumsi Pengembangan} \\
    Pengembangan media pembelajaran berbasis Augmented Reality ini didasarkan pada beberapa asumsi sebagai berikut:
    \begin{enumerate}
        \item Media pembelajaran berbasis Augmented Reality ini diharapkan dapat digunakan dalam pembelajaran matematika, khususnya pada materi bangun ruang sisi datar.
        \item Peserta didik akan lebih tertarik untuk belajar matematika ketika menggunakan media pembelajaran berbasis Augmented Reality.
        \item Penyajian bangun ruang sisi datar dalam bentuk virtual melalui teknologi AR dapat membantu peserta didik dalam memvisualisasikan konsep secara lebih nyata.
    \end{enumerate}

    \item \textbf{Keterbatasan Pengembangan} \\
    Penelitian pengembangan media pembelajaran berbasis Augmented Reality ini memiliki beberapa keterbatasan, antara lain:
    \begin{enumerate}
        \item Materi yang dikembangkan terbatas pada materi bangun ruang sisi datar.
        \item Subjek penelitian terbatas pada peserta didik kelas VIII SMP Negeri 3 Kalikajar.
        \item Media pembelajaran disusun berdasarkan Kurikulum 2013.
        \item Produk hanya dapat diakses secara online.
    \end{enumerate}
\end{enumerate}
